\documentclass[11pt]{article}

\usepackage[a4paper,margin=1in]{geometry}
\usepackage{amsmath,amssymb,bm}
\usepackage{graphicx}
\usepackage{booktabs}
\usepackage{siunitx}
\usepackage{subcaption}
\usepackage{xcolor}
\usepackage{hyperref}

\newcommand{\todo}[1]{\textcolor{red!70!black}{[TODO: #1]}}

\title{A molecularly resolved coarse-grained model of sarcomere assembly
under pulsatile activation}
\author{\todo{Author list}}
\date{\today}

\begin{document}

\maketitle

\begin{abstract}
Sarcomere assembly requires actin and myosin filaments to acquire a common
orientation and the axial registry that defines Z-discs and M-bands. This
organization develops in an active environment and is continually challenged
by contractile forces. The problem is especially relevant to hypertrophic
cardiomyopathy (HCM), in which diverse sarcomeric perturbations frequently
converge on excessive myosin-generated force, enhanced myofilament activation,
or incomplete relaxation. We develop a molecularly resolved coarse-grained
model that extends a finite-filament geometry with stochastic myosin
crossbridge cycling, discrete ACTN2 crosslinking, titin-like connections, and
local finite-valency myomesin coupling. Filament motion follows overdamped
Langevin dynamics, while molecular association and state changes are treated
kinetically. A pulsatile activation protocol permits direct comparison of
normal relaxation with elevated residual activation during diastole. The model
is designed to determine how the magnitude and timing of myosin activity alter
the formation and maintenance of sarcomeric order.
\end{abstract}


\section{Introduction}

Sarcomeres are built from interdigitating arrays of polar actin thin filaments
and bipolar myosin thick filaments. Their function depends not only on
filament alignment, but also on precise axial organization: actin barbed ends
meet at Z-discs, whereas thick-filament bare zones register at M-bands.
Although many of the molecular components are known, it remains unclear how
their local interactions produce ordered sarcomeres from a reorganizing
filament network.

Contractility makes this an intrinsically nonequilibrium assembly problem.
Myosin heads repeatedly attach to actin, generate force, and detach. These
forces can move filaments toward productive arrangements, but they can also
distort newly formed structures or trap incompatible contacts. Crosslinkers
and elastic scaffold proteins must therefore organize filaments while the
network remains mechanically active.

This balance may be altered in hypertrophic cardiomyopathy. HCM arises from
genetically and mechanistically diverse perturbations, but many disease-linked
changes in sarcomeric proteins converge on increased myosin activity,
heightened force production, or impaired relaxation. Incomplete relaxation is
particularly important because it leaves a larger fraction of the contractile
system active during diastole, when a normal sarcomere would experience a low
activation interval. Persistent diastolic activity could influence not only
instantaneous tension but also the longer-time assembly and remodeling of the
contractile lattice. % TODO: add HCM and diastolic dysfunction citations.

To examine this possibility, we build a molecular model on a coarse
finite-filament geometry. Individual myosin heads undergo stochastic binding
and working-stroke transitions, ACTN2 crosslinkers form discrete links between
thin filaments, and titin- and myomesin-like elements couple the Z-disc and
M-band regions. The model supports both tonic and heartbeat-like activation,
including a controlled residual activation level during diastole. Our goal is
to use this framework to ask how sustained myosin activity changes the
kinetics and spatial organization of sarcomere assembly, without prescribing
the final ordered structure.


\section{Molecular sarcomere model}

\subsection{Filament geometry and dynamics}

Actin and myosin are represented as finite, polar filaments with excluded
volume in a three-dimensional periodic domain. Each filament has a center
$\bm r_i$, an orientation $\bm u_i$, a length, and a radius. Myosin filaments
are bipolar, with two head-bearing halves separated by a central bare zone.

Filament positions and orientations evolve by overdamped Langevin dynamics.
In compact form,
\begin{align}
  d\bm r_i &= \bm M_i^{\mathrm t}\bm F_i\,dt
  + \sqrt{2k_{\mathrm B}T\bm M_i^{\mathrm t}}\,d\bm W_i^{\mathrm t},\\
  d\bm\theta_i &= \bm M_i^{\mathrm r}\bm\tau_i\,dt
  + \sqrt{2k_{\mathrm B}T\bm M_i^{\mathrm r}}\,d\bm W_i^{\mathrm r},
\end{align}
where $\bm F_i$ and $\bm\tau_i$ are the total force and torque and the mobility
tensors distinguish motion parallel and perpendicular to the filament axis.
Steric, motor, crosslink, titin, and myomesin interactions all contribute to
these mechanical updates.


\subsection{Actomyosin crossbridge cycle}

Myosin heads occupy fixed positions on the two halves of each thick filament.
A head can attach only to a nearby actin site and only the polarity-compatible
half of the bipolar filament generates a productive working stroke. Each head
cycles among detached, weakly bound, and strongly bound states. State changes
are sampled with kinetic Monte Carlo probabilities
\begin{equation}
  P(k,\Delta t)=1-\exp(-k\Delta t),
\end{equation}
where the relevant rate may depend on activation, elastic strain, and the
head's current chemical state.

A bound head behaves as an elastic connection between material sites on the
two filaments. The working stroke changes its axial rest position and thereby
generates force. The resulting forces act equally and oppositely on actin and
myosin, and off-center attachments also generate torques. Thick-filament
motion therefore emerges from the same crossbridge forces that move the thin
filaments rather than from a prescribed sliding velocity.


\subsection{Pulsatile activation and HCM-like perturbations}

The motor cycle is regulated by a periodic activation waveform $C(t)$ with a
high-activation interval and a low-activation diastolic interval. Local thin-
filament availability follows this input and can be reinforced by strongly
bound heads. Lower activation suppresses new head attachment, while heads
that are already bound complete their ordinary state-dependent turnover.

The reference protocol approaches a low activation floor during diastole. An
HCM-like protocol raises this residual activation while preserving the cycle
period and systolic waveform. This isolates the effect of incomplete
relaxation on assembly. Other disease-associated changes can be represented
separately by modifying head attachment, stroke, stiffness, or detachment
kinetics; these perturbations are not combined unless explicitly stated.


\subsection{ACTN2-mediated thin-filament crosslinking}

ACTN2 is represented by discrete crosslinking modules that bind nearby,
oppositely oriented actin filaments near their barbed-end regions. A newly
formed bridge is weak and can mature when it is both geometrically registered
and mechanically supported by productive actomyosin interactions. Weak and
mature bridges have distinct kinetics and mechanical responses.

Each module has two actin-binding arms. One arm may open while the other
remains attached to its material actin site, creating a transient tether. The
free arm can then bind a locally reachable site on the exposed portion of a
nearby actin or detach completely. ACTN2 attachment sites do not slide
continuously along actin. Finite contact between the ACTN2 connector and thick
filaments prevents a crosslink from passing freely through the actomyosin
overlap region.


\subsection{Titin- and myomesin-like connections}

Titin-like links connect a thick-filament half to a nearby actin filament that
has local Z-disc support. These links are tension bearing and persist across
the faster turnover of individual myosin heads. They provide longitudinal
coupling between thin- and thick-filament organization without fixing either
filament to an imposed spatial lattice.

Myomesin is represented by reversible local bridges between thick-filament
bare zones. Each bridge has an axial registry preference and a nonzero
transverse rest spacing. A finite number of binding sites limits the number of
simultaneous neighbors of each thick filament. Thus myomesin supplies local
M-region mechanics without introducing attraction to a prescribed global
M-band plane.


\subsection{Simulation protocol and observables}

Simulations begin from disordered actin positions and orientations, with
thick filaments weakly biased toward the longitudinal axis. A short
sterics-only equilibration removes severe overlaps before molecular
interactions are introduced. Interaction strengths are then increased
smoothly to avoid an impulsive mechanical start.

Assembly is quantified separately at the thin- and thick-filament levels.
Actin barbed-end density measures Z-disc organization, thick-filament bare-zone
density measures M-band organization, and filament centers report A-band
modulation. We additionally record filament orientation, crossbridge
occupancy, molecular state populations, forces, and the connectivity of the
ACTN2 and myomesin networks. Comparisons between normal and elevated
diastolic activation use matched initial conditions and identical mechanical
parameters.

Detailed interaction potentials, kinetic rates, numerical regularization,
checkpointing, and sensitivity analyses will be provided in the Supplementary
Methods.


% Results will be added after the molecular-model comparisons are complete.
% \section{Results}

\section*{Data and code availability}

\todo{Add repository, archived executable hashes, trajectories, and analysis
scripts used in the final study.}


% Enable once citations and bibliography entries are present:
% \bibliographystyle{unsrt}
% \bibliography{refs}

\end{document}
